\documentclass[12pt,a4paper]{article}
\usepackage[utf8]{inputenc}
\usepackage[T1]{fontenc}
\usepackage[spanish]{babel}
\usepackage{geometry}
\usepackage{booktabs}
\usepackage{longtable}
\usepackage{array}
\usepackage{xcolor}
\usepackage{listings}
\usepackage{hyperref}
\usepackage{fancyhdr}
\usepackage{graphicx}
\usepackage{titlesec}
\usepackage{enumitem}
\usepackage{mdframed}

\geometry{margin=2.5cm}

\definecolor{primaryblue}{RGB}{33,150,243}
\definecolor{criticalred}{RGB}{244,67,54}
\definecolor{normalgreen}{RGB}{76,175,80}
\definecolor{codegray}{RGB}{245,245,245}
\definecolor{codelisting}{RGB}{40,40,40}

\lstset{
  backgroundcolor=\color{codegray},
  basicstyle=\ttfamily\small,
  breaklines=true,
  frame=single,
  rulecolor=\color{gray},
  keywordstyle=\color{primaryblue}\bfseries,
  commentstyle=\color{normalgreen},
  stringstyle=\color{criticalred},
  showstringspaces=false,
  tabsize=2,
  captionpos=b,
}

\pagestyle{fancy}
\fancyhf{}
\rhead{EdgeLeak AI -- Actualización \#5}
\lhead{Documentación Técnica}
\cfoot{\thepage}

\hypersetup{
  colorlinks=true,
  linkcolor=primaryblue,
  urlcolor=primaryblue,
}

\title{
  \Huge\textbf{Documentación Técnica}\\[0.5em]
  \Large\textcolor{primaryblue}{EdgeLeak AI}\\[0.3em]
  \large Actualización \#5: Formulario de Registro, Validación de Correo\\
  y Corrección de Flujos de Navegación
}
\author{Proyecto EdgeLeak AI}
\date{\today}

\begin{document}

\maketitle
\thispagestyle{fancy}

\begin{mdframed}[linecolor=primaryblue,linewidth=1.5pt]
\textbf{Resumen Ejecutivo:} Esta actualización introduce cuatro mejoras críticas al sistema:
(1)~separación del campo \texttt{nombre} en cuatro campos independientes para mayor
precisión de datos; (2)~corrección de la validación de correo electrónico con regex;
(3)~corrección del flujo post-registro para evitar ingreso automático no autorizado;
(4)~corrección del flujo post-cambio de contraseña temporal para redirigir al login.
Se actualiza además el esquema de la base de datos local (versión 3, migración automática)
y el módulo CRUD de administración.
\end{mdframed}

\vspace{1em}
\tableofcontents
\newpage

% ============================================================
\section{Contexto y Motivación}
% ============================================================

Durante la fase de validación de usuario del sistema EdgeLeak AI se identificaron cuatro
defectos funcionales que comprometían la experiencia de usuario y la seguridad del sistema:

\begin{longtable}{>{\bfseries}p{3.5cm} p{4.5cm} p{5.5cm}}
\caption{Defectos identificados antes de la actualización} \label{tab:defectos} \\
\toprule
\textbf{ID} & \textbf{Descripción del Defecto} & \textbf{Impacto} \\
\midrule
\endfirsthead
\toprule
\textbf{ID} & \textbf{Descripción del Defecto} & \textbf{Impacto} \\
\midrule
\endhead
\midrule
\multicolumn{3}{r}{\small Continúa en la siguiente página\ldots} \\
\endfoot
\bottomrule
\endlastfoot
DEF-01 & Campo \texttt{nombre} único sin separación
         (primer nombre, segundo nombre, apellidos) &
         Imposibilidad de segmentar nombres para búsqueda, reportes
         y personalización correcta. \\[0.5em]
DEF-02 & Validación de correo electrónico ausente o ineficaz
         (sin regex aplicado) &
         El sistema aceptaba correos sin formato válido (p.\,ej.,
         sin símbolo \texttt{@}), causando registros inconsistentes. \\[0.5em]
DEF-03 & Flujo post-registro: al crear una cuenta, la navegación
         no garantizaba retorno al login &
         Posibilidad de acceso al panel sin autenticación explícita. \\[0.5em]
DEF-04 & Flujo post-cambio de contraseña temporal: el sistema
         redirigía directamente al \texttt{DashboardScreen} &
         Vulnerabilidad de seguridad; el usuario con contraseña
         temporal no realizaba una autenticación fresca con la
         nueva credencial. \\
\end{longtable}

% ============================================================
\section{Cambios Implementados}
% ============================================================

\subsection{DEF-01 -- Separación del Campo Nombre}

\subsubsection{Archivo afectado: \texttt{lib/data/models/usuario\_model.dart}}

El campo \texttt{nombre~(TEXT)} del modelo fue reemplazado por cuatro campos independientes.
Se añade un \emph{getter} calculado \texttt{nombre} para mantener compatibilidad hacia atrás
con toda la lógica existente de la UI.

\begin{longtable}{>{\ttfamily}p{4cm} p{2.5cm} p{7cm}}
\caption{Nuevos campos del modelo \texttt{UsuarioModel}} \label{tab:campos-modelo} \\
\toprule
\textbf{Campo (Dart)} & \textbf{Obligatorio} & \textbf{Descripción} \\
\midrule
\endfirsthead
\toprule
\textbf{Campo (Dart)} & \textbf{Obligatorio} & \textbf{Descripción} \\
\midrule
\endhead
\bottomrule
\endlastfoot
primerNombre   & Sí & Primer nombre del usuario. \\
segundoNombre  & No & Segundo nombre del usuario (valor por defecto: \texttt{''}). \\
primerApellido & Sí & Primer apellido del usuario. \\
segundoApellido& No & Segundo apellido del usuario (valor por defecto: \texttt{''}). \\
\end{longtable}

\textbf{Getter \texttt{nombre} (retrocompatibilidad):}
\begin{lstlisting}[language=Dart, caption={Getter calculado en UsuarioModel}]
String get nombre {
  final partes = <String>[primerNombre];
  if (segundoNombre.isNotEmpty) partes.add(segundoNombre);
  partes.add(primerApellido);
  if (segundoApellido.isNotEmpty) partes.add(segundoApellido);
  return partes.join(' ');
}
\end{lstlisting}

\textbf{Lógica de \texttt{fromMap} con compatibilidad de migración:}
\begin{lstlisting}[language=Dart, caption={Deserialización con soporte a datos legados}]
factory UsuarioModel.fromMap(Map<String, dynamic> map) {
  final primerNombreLegado = map['nombre'] as String? ?? '';
  final primerNombreNuevo  = map['primer_nombre'] as String? ?? '';
  return UsuarioModel(
    id: map['id'],
    primerNombre:    primerNombreNuevo.isNotEmpty
                       ? primerNombreNuevo : primerNombreLegado,
    segundoNombre:   map['segundo_nombre']  as String? ?? '',
    primerApellido:  map['primer_apellido'] as String? ?? '',
    segundoApellido: map['segundo_apellido']as String? ?? '',
    correo:    map['correo'],
    password:  map['password'],
    esTemporal: map['es_temporal'] ?? 0,
    rol:        map['rol'] ?? 'operador',
  );
}
\end{lstlisting}

\subsubsection{Archivo afectado: \texttt{lib/data/services/database\_service.dart}}

El esquema SQLite se actualizó a \textbf{versión 3}. Se implementó migración automática
\texttt{onUpgrade} para instalaciones existentes (versión 2 $\rightarrow$ 3).

\begin{longtable}{>{\ttfamily}p{4.5cm} p{3.5cm} p{2.5cm} p{2.5cm}}
\caption{Nuevas columnas en la tabla \texttt{usuarios} (v3)} \label{tab:columnas-db} \\
\toprule
\textbf{Columna SQL} & \textbf{Tipo} & \textbf{Default} & \textbf{NULL?} \\
\midrule
\endfirsthead
\toprule
\textbf{Columna SQL} & \textbf{Tipo} & \textbf{Default} & \textbf{NULL?} \\
\midrule
\endhead
\bottomrule
\endlastfoot
primer\_nombre  & TEXT & \texttt{''} & No \\
segundo\_nombre & TEXT & \texttt{''} & Sí \\
primer\_apellido& TEXT & \texttt{''} & No \\
segundo\_apellido& TEXT & \texttt{''} & Sí \\
\end{longtable}

\textbf{Script de migración (onUpgrade, oldVersion $<$ 3):}
\begin{lstlisting}[language=SQL, caption={Migración v2 -> v3 en SQLite}]
ALTER TABLE usuarios ADD COLUMN primer_nombre  TEXT DEFAULT '';
ALTER TABLE usuarios ADD COLUMN segundo_nombre TEXT DEFAULT '';
ALTER TABLE usuarios ADD COLUMN primer_apellido  TEXT DEFAULT '';
ALTER TABLE usuarios ADD COLUMN segundo_apellido TEXT DEFAULT '';
-- Preservar datos legados: migrar 'nombre' a 'primer_nombre'
UPDATE usuarios
  SET primer_nombre = nombre
  WHERE nombre IS NOT NULL AND nombre != '';
\end{lstlisting}

\subsubsection{Archivo afectado: \texttt{lib/controllers/auth\_controller.dart}}

La firma del método \texttt{register()} se actualizó para recibir los cuatro campos de nombre
como parámetros posicionales independientes:

\begin{lstlisting}[language=Dart, caption={Nueva firma de AuthController.register()}]
Future<bool> register(
  String primerNombre,
  String segundoNombre,
  String primerApellido,
  String segundoApellido,
  String correo,
  String password, {
  String rol = 'operador',
}) async { ... }
\end{lstlisting}

\subsubsection{Archivos afectados: \texttt{register\_screen.dart} y \texttt{admin\_users\_screen.dart}}

El formulario de registro fue rediseñado con cuatro campos de texto individuales.
El diálogo de edición en el panel CRUD también se actualizó con los cuatro campos.

\begin{longtable}{p{5cm} p{3cm} p{5.5cm}}
\caption{Campos del formulario de registro (nuevo diseño)} \label{tab:form-registro} \\
\toprule
\textbf{Campo en UI} & \textbf{Obligatorio} & \textbf{Validación} \\
\midrule
\endfirsthead
\toprule
\textbf{Campo en UI} & \textbf{Obligatorio} & \textbf{Validación} \\
\midrule
\endhead
\bottomrule
\endlastfoot
Primer Nombre *    & Sí & No vacío (validado al pulsar REGISTRARSE). \\
Segundo Nombre     & No & Ninguna. \\
Primer Apellido *  & Sí & No vacío (validado al pulsar REGISTRARSE). \\
Segundo Apellido   & No & Ninguna. \\
Correo Electrónico *& Sí & Regex de correo válido (ver §2.2). \\
Contraseña *       & Sí & Política de contraseñas fuerte (4 criterios). \\
Código Admin       & No & Comparado contra variable de entorno \texttt{SECRET\_ADMIN\_CODE}. \\
\end{longtable}

% ------------------------------------------------------------
\subsection{DEF-02 -- Validación de Correo Electrónico con Regex}
% ------------------------------------------------------------

\subsubsection{Problema}

Anteriormente, el campo de correo únicamente tenía configurado el tipo de teclado
\texttt{TextInputType.emailAddress}, lo que no impide en modo desktop/web o en dispositivos
físicos la escritura de cadenas sin el símbolo \texttt{@} ni dominio.

\subsubsection{Solución}

Se añadió un objeto \texttt{RegExp} estático a \texttt{\_RegisterScreenState} con el patrón
estándar RFC-5322 simplificado:

\begin{lstlisting}[language=Dart, caption={Regex de validación de correo en RegisterScreen}]
static final _emailRegex = RegExp(
  r'^[a-zA-Z0-9._%+\-]+@[a-zA-Z0-9.\-]+\.[a-zA-Z]{2,}$',
);

bool _validarCorreo(String correo) {
  return _emailRegex.hasMatch(correo);
}
\end{lstlisting}

La validación se ejecuta en dos momentos:

\begin{itemize}[leftmargin=2em]
  \item \textbf{En tiempo real} (\texttt{onChanged}): muestra el mensaje de error directamente
        en el campo mediante la propiedad \texttt{errorText} del \texttt{InputDecoration}.
  \item \textbf{Al pulsar REGISTRARSE}: bloquea el proceso y muestra el error si el formato
        es inválido o el campo está vacío.
\end{itemize}

\begin{longtable}{p{5cm} p{3.5cm} p{5cm}}
\caption{Ejemplos de validación de correo} \label{tab:email-validacion} \\
\toprule
\textbf{Cadena de entrada} & \textbf{Resultado} & \textbf{Motivo} \\
\midrule
\endfirsthead
\toprule
\textbf{Cadena de entrada} & \textbf{Resultado} & \textbf{Motivo} \\
\midrule
\endhead
\bottomrule
\endlastfoot
\texttt{usuario@dominio.com}  & \textcolor{normalgreen}{\textbf{Válido}}   & Formato RFC completo. \\
\texttt{user.name+tag@sub.d.co}& \textcolor{normalgreen}{\textbf{Válido}}  & Caracteres especiales permitidos. \\
\texttt{sinArroba.com}        & \textcolor{criticalred}{\textbf{Inválido}} & Ausencia de símbolo \texttt{@}. \\
\texttt{@sinusuario.com}      & \textcolor{criticalred}{\textbf{Inválido}} & Parte local vacía. \\
\texttt{usuario@}             & \textcolor{criticalred}{\textbf{Inválido}} & Dominio ausente. \\
\texttt{usuario@dom}          & \textcolor{criticalred}{\textbf{Inválido}} & TLD $<$ 2 caracteres. \\
\end{longtable}

% ------------------------------------------------------------
\subsection{DEF-03 -- Corrección del Flujo Post-Registro}
% ------------------------------------------------------------

\subsubsection{Problema}

El método de navegación \texttt{Navigator.pop(context)} usado tras un registro exitoso
puede comportarse de forma impredecible cuando el stack de navegación no contiene la pantalla
de login como inmediato predecesor (por ejemplo, en escenarios de deep-linking o reload
de estado). Esto podría permitir acceso sin autenticación.

\subsubsection{Solución}

Se reemplazó \texttt{Navigator.pop(context)} por
\texttt{Navigator.pushNamedAndRemoveUntil} que limpia el stack de navegación completo
y establece el login como única pantalla activa:

\begin{lstlisting}[language=Dart, caption={Navegación post-registro corregida}]
// ANTES (riesgoso)
Navigator.pop(context);

// AHORA (seguro)
Navigator.pushNamedAndRemoveUntil(
  context,
  AppRoutes.loginScreen,
  (route) => false,   // elimina todas las rutas previas
);
\end{lstlisting}

% ------------------------------------------------------------
\subsection{DEF-04 -- Corrección del Flujo Post-Cambio de Contraseña Temporal}
% ------------------------------------------------------------

\subsubsection{Problema}

Al cambiar la contraseña obligatoria (flujo de contraseña temporal), el sistema ejecutaba:
\begin{lstlisting}[language=Dart]
Navigator.pushReplacementNamed(context, AppRoutes.dashboardScreen);
\end{lstlisting}
Esto ingresaba directamente al panel sin requerir una nueva autenticación con la nueva
contraseña, vulnerando el principio de ``sesión limpia post-credencial''.

\subsubsection{Solución}

Se realizaron dos cambios coordinados:

\textbf{1. En \texttt{AuthController.cambiarPasswordObligatorio()}:}
\begin{lstlisting}[language=Dart, caption={Limpieza de sesión al cambiar contraseña}]
// ANTES: actualizaba usuarioActual con nueva contraseña
usuarioActual = UsuarioModel(id: ..., password: nuevaPassword, ...);

// AHORA: limpia la sesión para forzar re-autenticación
usuarioActual = null;
\end{lstlisting}

\textbf{2. En \texttt{ChangePasswordScreen} (botón GUARDAR Y CONTINUAR):}
\begin{lstlisting}[language=Dart, caption={Redirección a login tras cambio exitoso}]
// ANTES
Navigator.pushReplacementNamed(context, AppRoutes.dashboardScreen);

// AHORA
Navigator.pushNamedAndRemoveUntil(
  context,
  AppRoutes.loginScreen,
  (route) => false,
);
\end{lstlisting}

El mensaje de confirmación también fue actualizado para ser informativo:
\begin{quote}
\textit{``Contraseña actualizada. Inicia sesión con tu nueva clave.''}
\end{quote}

% ============================================================
\section{Diagrama de Flujo de Navegación (Corregido)}
% ============================================================

\begin{longtable}{p{5cm} p{5cm} p{3.5cm}}
\caption{Flujo de navegación antes y después de la actualización} \label{tab:flujo-nav} \\
\toprule
\textbf{Acción del usuario} & \textbf{Pantalla origen} & \textbf{Destino correcto} \\
\midrule
\endfirsthead
\toprule
\textbf{Acción del usuario} & \textbf{Pantalla origen} & \textbf{Destino correcto} \\
\midrule
\endhead
\bottomrule
\endlastfoot
Registro exitoso & \texttt{RegisterScreen} & \texttt{LoginScreen} (stack limpio) \\
Login con clave temporal & \texttt{LoginScreen} & \texttt{ChangePasswordScreen} \\
Cambio de clave exitoso & \texttt{ChangePasswordScreen} & \texttt{LoginScreen} (stack limpio) \\
Login normal & \texttt{LoginScreen} & \texttt{DashboardScreen} \\
\end{longtable}

% ============================================================
\section{Actualización de Pruebas Unitarias (TDD)}
% ============================================================

El archivo \texttt{test/unit\_tests/user\_logic\_test.dart} fue actualizado para cubrir
los nuevos campos de nombre y los nuevos comportamientos del modelo:

\begin{longtable}{p{6cm} p{7.5cm}}
\caption{Nuevos casos de prueba añadidos} \label{tab:tests} \\
\toprule
\textbf{Nombre del Test} & \textbf{Propósito} \\
\midrule
\endfirsthead
\toprule
\textbf{Nombre del Test} & \textbf{Propósito} \\
\midrule
\endhead
\bottomrule
\endlastfoot
\textit{getter nombre combina todos los campos} &
  Verifica que \texttt{nombre} devuelve la concatenación correcta de los 4 campos. \\[0.5em]
\textit{getter nombre omite campos opcionales vacíos} &
  Valida que \texttt{segundoNombre} y \texttt{segundoApellido} vacíos no generen espacios dobles. \\[0.5em]
\textit{toMap incluye los cuatro campos de nombre} &
  Garantiza que la serialización escribe \texttt{primer\_nombre}, \texttt{segundo\_nombre}, etc. \\[0.5em]
\textit{fromMap reconstruye el modelo correctamente} &
  Verifica que la deserialización desde nuevos campos funciona. \\[0.5em]
\textit{fromMap usa el campo nombre legado} &
  Garantiza retrocompatibilidad: si \texttt{primer\_nombre} está vacío, usa el campo \texttt{nombre} legado. \\
\end{longtable}

% ============================================================
\section{Resumen de Archivos Modificados}
% ============================================================

\begin{longtable}{p{7cm} p{6.5cm}}
\caption{Resumen de todos los archivos modificados en esta actualización} \label{tab:archivos} \\
\toprule
\textbf{Archivo} & \textbf{Cambio Principal} \\
\midrule
\endfirsthead
\toprule
\textbf{Archivo} & \textbf{Cambio Principal} \\
\midrule
\endhead
\bottomrule
\endlastfoot
\texttt{lib/data/models/usuario\_model.dart} &
  4 campos de nombre, getter \texttt{nombre}, \texttt{fromMap} con fallback legado. \\[0.5em]
\texttt{lib/data/services/database\_service.dart} &
  Versión 3, nuevas columnas, migración automática. \\[0.5em]
\texttt{lib/controllers/auth\_controller.dart} &
  \texttt{register()} con 4 params; \texttt{cambiarPasswordObligatorio()} limpia sesión. \\[0.5em]
\texttt{lib/view/screens/register\_screen.dart} &
  4 campos de nombre, validación regex de correo, navegación segura al login. \\[0.5em]
\texttt{lib/view/screens/change\_password\_screen.dart} &
  Redirección al login con stack limpio tras cambio exitoso. \\[0.5em]
\texttt{lib/view/screens/admin\_users\_screen.dart} &
  Diálogo de edición actualizado con 4 campos de nombre. \\[0.5em]
\texttt{test/unit\_tests/user\_logic\_test.dart} &
  5 nuevos tests; actualización de tests existentes con nueva API. \\
\end{longtable}

% ============================================================
\section{Consideraciones de Seguridad}
% ============================================================

\begin{itemize}[leftmargin=2em]
  \item \textbf{Principio de sesión limpia:} Al limpiar \texttt{usuarioActual = null} en
        \texttt{cambiarPasswordObligatorio()}, el token de sesión en memoria queda invalidado.
        Cualquier pantalla que dependa de \texttt{authController.usuarioActual?.esAdmin}
        evaluará correctamente a \texttt{false}, impidiendo acceso accidental al CRUD
        de administración.

  \item \textbf{Stack de navegación limpio:} El uso de \texttt{pushNamedAndRemoveUntil} con
        predicado \texttt{(route) => false} garantiza que no existan rutas residuales en el
        historial que permitan navegar hacia atrás al dashboard sin credenciales.

  \item \textbf{Validación en cliente:} La validación de correo con regex es una capa de
        defensa de UX. En una fase futura, se recomienda añadir validación en el servicio
        de base de datos (restricción \texttt{CHECK} en SQLite o validación en el backend
        si se migra a servidor remoto).
\end{itemize}

\end{document}
